\subsection{What the seven statements actually say}
\paragraph{J0 (Textual correction and scope).}
We keep the uncertain exponent in the transcription of Theorem~3
at 44:2 uncertain.  The letter at 44:8 unambiguously writes
``$J^i$ et $J^{n+1-i}$ sont canoniquement duales l'une de l'autre'',
so $n+1-i$ is the mathematically intended complementary index.
This does not resolve an illegible handwritten character by fiat.
The same leaf says that the equivalence relation defining $J^i$ is
one ``que je ne sais pas bien identifier pour l'instant''.
Consequently the literal quotient by the manuscript's Albanese
equivalence is not automatically identified with Murre's modern
universal regular quotient.  That identification needs proof.

For reference, the seven assertions separate as follows.  This table
records present proof status, rather than representing unresolved
assertions as completed items.
\begin{description}
\item[Theorem 1, 44:2.]
{\raggedright The Codex transcription reads ``$N(\zeta)\subset\operatorname{Alb}(S)$
n'est autre \emph{ensemblistement} que le noyau'', while Pro reads
\emph{en substance}; this conflict is retained. The common
incidence-kernel interpretation is proved in J1i below, even
scheme-theoretically. The common kernel for all regular tests also
has the stated set-theoretic property. Identifying either with the
particular undefined manuscript $N(\zeta)$ remains a separate
textual and mathematical gate.\par}
\item[Theorem 2, 44:2.]
The quotient dimensions are bounded, ``en fait, par
$\frac12b^{2i-1}(X)$''.  For a modern algebraic representative in
codimensions $1,2,n$ this bound follows by taking a miniversal
correspondence and its injective rational Tate realization; see
\cite[Proposition 8.2, p.~31]{DGJ_ACMVFunctor}.  Specifically choose
a prime different from the characteristic and the degree of the
miniversal correspondence; the injection has source of rank
$2\dim\operatorname{Ab}^i(X)$ and target of rank $b^{2i-1}$.
For incidence equivalence, Saito's Proposition~3.4 gives this bound
in every codimension; J1i proves the exact family-quotient form.
For the equivalence defined by all regular homomorphisms, the
general uniform-boundedness problem remains unresolved.
\item[Theorem 3, 44:2.]
The quotient is naturally an abelian variety, with complementary
duality.  The incidence quotient exists in every codimension and has
complementary isogeny duality by Saito's Theorem~3.6 and
Proposition~4.4. Existence of the modern universal regular
representative is known in codimensions $1,2,n$; its general
existence, identification with the incidence quotient, and the
claimed integral duality are separate issues.  Under the modern
algebraic-representative interpretation explicitly requested in
J.2, the literal integral duality is false by J2a below, already
in characteristic five for a three-dimensional Jacobian.  This
does not identify the manuscript's less precisely specified
incidence quotient.  General regular representability and that identification remain
unresolved; J1i gives the exact established incidence version and
its specified isogeny normalization.
\item[Theorem 4, 44:2.]
The raw middle self-duality is ``associée à une polarisation''.
This is false, even allowing either overall sign: J2b below.
The primitive qualification in the margin of 44:8 is essential.
There is also an unresolved margin-reading conflict at 44:2:
Codex reads a form of ``Annulé/Cancelé tel quel'', whereas Pro
reads ``Connu tel quel''; this is not silently resolved.
\item[Theorem 5, 44:4.]
The Néron--Tate pairing is $(-1)^mZ\cdot Z'$ for cycles of
codimensions $i,j$ with $i+j=n$, whose restrictions to the generic
member of a hyperplane pencil are algebraically trivial, with all
pencil fibres geometrically integral. This literal formula is false:
J5 below gives a split quadric fourfold in characteristic three and
four explicit planes with both generic restrictions literally zero,
while their signed intersection has degree two. Here $i=j=m=2$,
so the retained sign variable causes no ambiguity. The example
also meets every explicit integral-fibre hypothesis. A stronger
height comparison with admissibility or vertical-correction
conditions is a different statement and is not refuted.
\item[Theorem 6, 44:4.]
The Codex transcription reads \emph{isogénie}, whereas the Pro
transcription reads \emph{isomorphisme}, if $2i\leq n-1$;
without the scan this disagreement must remain explicit.  In
either reading, its image
is the fixed factor if $2i=n$; it is onto if $2i\geq n+1$.
These are assertions about algebraic representatives, not merely
cohomological weak Lefschetz. The literal middle clause is false:
J6C gives a smooth cubic fourfold and a pencil with all fibres
geometrically integral whose target has a constant elliptic factor,
although the source representative is zero. The axis is singular,
which the manuscript does not exclude. This works for both the
incidence and universal regular quotients. J6P proves the
codimension-two fixed-image clause for a smooth pencil axis, and
generic stable-range surjectivity in every characteristic; the
finite-kernel step retains an explicit torsion hypothesis. J6Z
proves surjectivity in the highest nonzero target codimension.
J6L below proves
the stable isogeny under the relevant algebraic Lefschetz lowering
operator on $X$, and surjectivity under the corresponding operator
on $Y$, with no claim that those hypotheses are automatic.
\item[Theorem 7, 44:6.]
``L'équivalence numérique $=$ $\tau$-équivalence'' for all cycles
on a smooth projective variety.  It is false in every intermediate codimension in which such
a counterexample can occur: the explicit Fermat-quartic Ceresa
cycle below is numerically trivial but is not $\tau$-equivalent
to zero.
\end{description}

\subsection{The precise incidence quotient and the remaining regular-quotient problem}
\paragraph{J1i (Known construction, with the family-kernel consequence proved here).}
Let $k$ be an algebraically closed field of any characteristic, let
$X/k$ be smooth projective of pure dimension $n$, and write
$A^i(X)=\operatorname{CH}^i(X)_{\rm alg}$. Define
\[
 I^i(X)=\bigcap_{T,w}\ker\left(
 w_*:A^i(X)\longrightarrow\operatorname{Pic}^0(T)(k)\right),
 \qquad w\in\operatorname{CH}^{n-i+1}(X\times T),
\]
where $T$ runs through smooth projective $k$-varieties; on $k$-points
$\operatorname{Pic}^0(T)$ means its reduced identity component.
This is \emph{incidence equivalence}, also called Picard equivalence
in the older sources. It is not a definition by all regular
homomorphisms to abelian varieties.

Kleiman explicitly separates the two relations in
\cite[pp.~446--447]{DGJ_Kleiman1971}. His p.~446 attributes the
boundedness argument for incidence equivalence to an unpublished
1964 observation of Grothendieck, and p.~447, footnote~5, records
the complementary duality with finite kernel. Thus the historical
connection with 44:2--9 has a primary-source basis, without
identifying two different equivalence relations.
The complete construction is Saito's
\cite[Definition 3.1, p.~108; Proposition 3.4, pp.~109--110;
Theorem 3.6, pp.~110--111]{DGJ_Saito1979}: there is an abelian
variety $P^i_{\rm inc}(X)$ and a surjective regular Picard
homomorphism
\[
 \theta_X^i:A^i(X)\longrightarrow P^i_{\rm inc}(X)(k),
 \qquad\ker\theta_X^i=I^i(X),
 \qquad 2\dim P^i_{\rm inc}(X)\leq b^{2i-1}(X).
\]
Here the Betti number is the dimension of $\ell$-adic cohomology
for any prime $\ell\ne\operatorname{char}k$. The map is universal
among Picard homomorphisms, which are the homomorphisms obtained
by a complementary correspondence to the Picard variety of a
smooth projective variety, with the image viewed as an abelian
subvariety. All these statements hold in every codimension,
including characteristics two and three. Saito's bound is written
using $b^{2n-2i+1}$; Poincaré duality gives the displayed index.

Here is the exact family-kernel statement which follows. Let
$(S,s_0)$ be smooth projective and pointed, and let
$\xi\in\operatorname{CH}^i(S\times X)$. Correspondence
functoriality supplies a homomorphism of $k$-abelian varieties
\[
 f_\xi:\operatorname{Alb}(S)\longrightarrow P^i_{\rm inc}(X),
 \qquad
 f_\xi([s]-[s_0])=\theta_X^i(\xi_s-\xi_{s_0}).
\]
For each $w$ in the definition of $I^i(X)$, composition with
$\xi$ is a divisor correspondence on $S\times T$ and defines
$g_{w,\xi}:\operatorname{Alb}(S)\to\operatorname{Pic}^{\rm var}(T)$.
Define $N_{\rm inc}(\xi)$ as the scheme-theoretic intersection of
all their kernels. Then
\[
 N_{\rm inc}(\xi)=\ker f_\xi,
 \qquad
 \operatorname{Alb}(S)/N_{\rm inc}(\xi)
     \simeq\operatorname{Im}(f_\xi),
 \qquad
 \dim\operatorname{Alb}(S)/N_{\rm inc}(\xi)
       \leq\tfrac12 b^{2i-1}(X).                 \tag{J1i.1}
\]
The quotient and the image here are abelian varieties; the kernel
may be nonreduced. To prove the equality, universality makes every
$g_{w,\xi}$ factor through $f_\xi$, giving one inclusion.
Conversely, $\theta_X^i$ is itself a Picard homomorphism
\cite[(3.5), p.~110]{DGJ_Saito1979}. Thus for some $T_0,w_0$,
its composition with a closed immersion
$P^i_{\rm inc}(X)\hookrightarrow\operatorname{Pic}^{\rm var}(T_0)$
is $(w_0)_*$. The corresponding $g_{w_0,\xi}$ has kernel exactly
$\ker f_\xi$, giving the other inclusion as group schemes, not
only on geometric points. The fppf quotient by a homomorphism's
kernel is its abelian-variety image. The dimension assertion now
follows from Saito's bound. This proves (J1i.1).

Saito further constructs, using his explicitly specified
miniversal-cycle integers, an isogeny
\[
 P^i_{\rm inc}(X)\longrightarrow
          P^{n+1-i}_{\rm inc}(X)^\vee;
\]
see \cite[(4.1)--(4.2), p.~112; Proposition 4.4 and
Corollary 4.5, p.~113]{DGJ_Saito1979}. This is a known
\emph{isogeny} statement. Its integer factors must be retained;
it does not supply the unscaled integral intersection isomorphism
asserted in J.2. Therefore the incidence reading of the
representability and boundedness portions of Theorems~1--3 is
settled by the cited construction and (J1i.1). The stronger
integral duality and the identification with a quotient defined
by all regular maps are separate assertions.

\paragraph{The regular interpretation and its exact remaining gate.}
A homomorphism $h:A^i(X)\to B(k)$ to a $k$-abelian variety is
\emph{regular} when every pointed smooth projective family of
cycles induces a morphism to $B$. Put
\[
 I^i_{\rm ab}(X)=\bigcap_h\ker h.
\]
This is Samuel's abelian equivalence. Every Picard homomorphism is
regular, so
$I^i_{\rm ab}(X)\subseteq I^i(X)$. A universal regular
homomorphism, if it exists, therefore has a canonical surjective
homomorphism from its abelian variety to $P^i_{\rm inc}(X)$.
This observation does not prove that the kernel is finite or zero.

For a fixed family $\xi$, let $N_{\rm ab}(\xi)$ be the
scheme-theoretic intersection of the kernels of all homomorphisms
$h_\xi:\operatorname{Alb}(S)\to B$ obtained from regular $h$.
Then the induced group homomorphism
\[
 \operatorname{Alb}(S)(k)\longrightarrow
                    A^i(X)/I^i_{\rm ab}(X)
\]
has kernel $N_{\rm ab}(\xi)(k)$, without assuming a universal
regular representative. Indeed, every composite on degree-zero
zero-cycles of $S$ is regular and factors through its Albanese;
hence the Albanese kernel is killed modulo $I^i_{\rm ab}(X)$.
The displayed map is well defined because the Albanese map on
zero-cycles is surjective over an algebraically closed field, and
its kernel is the stated intersection by definition. The
intersection is a closed subgroup scheme: on a noetherian
scheme the sum of the defining ideal sheaves is generated by a
finite subfamily. Thus each individual family has an abelian
quotient. What is missing is a uniform bound as the family varies.

More precisely, the existence of such a uniform dimension bound
is equivalent to existence of a universal regular homomorphism.
The bounded-target criterion is
\cite[Theorem 2.2(ii), p.~103 and proof pp.~103--106]{DGJ_Saito1979}
for the class of all regular homomorphisms, identified in his
Example (1.3.1), p.~101. To apply it to the family formulation,
any surjective regular $h$ admits a family on an abelian variety
whose induced map to its target is surjective
\cite[Proposition 1.2, pp.~98--100]{DGJ_Saito1979}; this target
is a quotient of the corresponding family quotient, so a uniform
family bound bounds all regular targets. Conversely, a universal
regular representative is itself among the tests defining every
$N_{\rm ab}(\xi)$, so each such quotient is its abelian
subvariety image. This proves the claimed equivalence of the two
boundedness problems. A bound by the particular Betti number is
stronger than mere existence.

The known existence cases are $i=1,2,n$.
For codimension two the complete repair of Murre's proof is
\cite[Theorem 1 and its two proofs, author version
pp.~2--4]{DGJ_Kahn2021}; for all three indices the Betti bound is
\cite[Proposition 8.2, p.~31]{DGJ_ACMVFunctor}.
That proposition explicitly restricts its cohomological injection
to $i=1,2,n$. Applying it in arbitrary codimension would assume
the missing injectivity of the higher Bloch map and would falsely
close the remaining gate.

A concrete unresolved test is Kahn's Conjecture~1, author version
p.~5: take a generic abelian threefold $A$ defined over a field
$k_0$ finitely generated over $\mathbb Q$, a nonisotrivial elliptic
curve $E/k_0(t)$, and
\[
 V=(A_{k_0(t)}\times E)_{\overline{k_0(t)}}.
\]
The conjecture states that a universal regular homomorphism on
$A^3(V)$ does not exist. Its preceding infinite-torsion result
alone does not prove this conjecture. Although every surjective
regular homomorphism is surjective on prime-to-characteristic
primary torsion \cite[Remark 2, p.~4]{DGJ_Kahn2021}, an infinite
torsion kernel may be killed in every finite-dimensional target.
To refute universal existence one must still construct regular
quotients of unbounded dimension, or another actual obstruction
satisfying the regularity condition. Neither the torsion argument
nor the opened motivic constructions provides that step here.
Accordingly the all-regular version in codimension three on
fourfolds, and its requested general form, remain unresolved.

\paragraph{A single self-incidence kernel is not the common kernel.}
The letter at 44:9 constructs $N$ from one symmetric
self-correspondence on a parameter variety. Such a kernel cannot
be substituted for $N_{\rm inc}(\xi)$ without an additional
argument. The distinction is visible in the same explicit
threefold as J2b. Let $X=E\times\mathbb P^1\times\mathbb P^1$,
let $S=E^\vee$, and set $\xi=c_1(\mathcal P)h_1$, where
$\mathcal P$ is the normalized Poincaré bundle. Its
self-incidence divisor is zero because $h_1^2=0$, so its
self-kernel is all of $S$. But pairing with the independent
family $c_1(\mathcal P)h_2$ gives the elliptic duality
$S\to S^\vee$, up to sign, which is an isomorphism. Hence
$N_{\rm inc}(\xi)=0$. Equivalently the family map is the
inclusion of the first elliptic factor in $P^2_{\rm inc}(X)=S^2$.
The projective-bundle decomposition and the two complementary
Poincaré tests prove this identification integrally, just as in
J2b. Thus one self-pairing can miss an entire abelian family that
other incidence tests detect. This refutes that particular
shortcut from 44:9 to Theorem~1; it does not decide an undefined
$N(\xi)$ on 44:2 by assigning it a new definition.

\paragraph{Exact status of the remaining Lefschetz assertions.}
The literal middle clause of Theorem~6 is refuted by J6C below.
The stable- and high-range clauses remain unresolved in their
full stated generality. Even after choosing the incidence quotient,
they concern algebraic correspondences and abelian varieties.
Cohomological weak Lefschetz identifies the ambient
cohomology in its stable range, but does not show that an inverse
on the required algebraic part is induced by correspondences.
For the universal regular quotient, general existence is already
the preceding gate. For the literal manuscript, the two
transcriptions additionally disagree between an isomorphism and
an isogeny in the stable range. Neither a cohomological
isomorphism nor Saito's complementary isogeny proves the stronger
integral restriction claim. These are retained as explicit
unresolved assertions. J6L, J6P, and J6Z establish the precisely
stated conditional and restricted cases below.

\subsection{An explicit counterexample to raw positivity}
\paragraph{J2b (Refuted, with a complete counterexample).}
Let $k=\overline{\mathbb F}_3$ and let
\[
 E:\ y^2z=x^3-xz^2\subset\mathbb P^2_k,
 \qquad X=E\times\mathbb P^1\times\mathbb P^1.
\]
The curve is nonsingular: in characteristic three the derivative
of $y^2z-x^3+xz^2$ with respect to $x$ is $z^2$; if all partials
vanished, $z=0$, the equation would force $x=0$, and the derivative
with respect to $z$ would then force $y=0$, which is not a
projective point.  Its point at infinity is the origin.  Thus $X$
is a smooth projective threefold satisfying exactly the hypotheses
of J.2.  In fact the same construction works with any elliptic
curve in every characteristic.

Let $h_1,h_2$ be the two hyperplane classes.  The projective bundle
formula for Chow groups \cite[Tag 02TX]{DGJ_StacksPB} gives
\[
 \operatorname{CH}^2(X)=
   \operatorname{Pic}(E)h_1\oplus\operatorname{Pic}(E)h_2
          \oplus\mathbb Z h_1h_2.
\]
Consequently its algebraically trivial subgroup is
$\operatorname{Pic}^0(E)h_1\oplus\operatorname{Pic}^0(E)h_2$.
The decomposition and its inverse are given by multiplication by
the hyperplane classes and proper pushforward along the
projections.  They apply to families and preserve regular
homomorphisms.  The universal property of the Picard variety
therefore proves, as an algebraic representative and not just as
an abstract group of points,
\[
 \operatorname{Ab}^2(X)=E^\vee\times E^\vee\simeq E^2.
\]
The last identification uses the standard principal polarization
of the elliptic curve.

Write $T=(E^\vee)^2$, with origin $a=0$, and let
$\mathcal P_1,\mathcal P_2$ be the pullbacks to $T\times E$ of
the normalized Poincaré bundle on $E^\vee\times E$.
The completely specified family
\[
 z=c_1(\mathcal P_1)h_1+c_1(\mathcal P_2)h_2
       \in\operatorname{CH}^2(T\times X)
\]
has $z(0)=0$ and induces the identity on the displayed algebraic
representative.  Form exactly the divisor correspondence in the
letter at 44:9:
\[
 D=r_*(p^*z\cdot q^*z)\in\operatorname{CH}^1(T\times T).
\]
Because $h_1^2=h_2^2=0$ and $\int_{\mathbb P^1\times\mathbb P^1}
h_1h_2=1$, the two diagonal entries of its homomorphism
$\rho:T\to T^\vee$ vanish and the two off-diagonal entries are
the elliptic intersection duality.  Thus, up to one overall sign
depending on the Poincaré convention,
\[
 \rho=\begin{pmatrix}0&\lambda_E\\\lambda_E&0\end{pmatrix}.
\]
There is no hidden scalar two in either entry: the two cross terms
give the two distinct off-diagonal entries, once each.
This homomorphism is an isomorphism, with inverse obtained by
interchanging the two dual factors using $\lambda_E^{-1}$.

Pull it back along the diagonal $u\mapsto(u,u)$ and the
antidiagonal $u\mapsto(u,-u)$ of $E^2$.  The resulting symmetric
homomorphisms are respectively $2\lambda_E$ and $-2\lambda_E$,
with both signs reversed if $-\rho$ is used.  The restriction of
an ample line bundle to either elliptic subgroup must have positive
degree.  One of these two degrees is negative for either overall
sign.  Hence neither $\rho$ nor $-\rho$ is a polarization.
The same argument rules out positivity of a rational
Néron--Severi class inducing this correspondence.  Since $\rho$
has zero kernel, passing to the letter's quotient $T/\ker\rho$
does not change the counterexample.

This refutes the literal unqualified Theorem~4 and the positivity
assertion in J.2, including the proposed single sign from
Theorem~5.  It does not refute the primitive correction which
Grothendieck himself adds at 44:8:
``il faut tenir compte de la décomposition primitive''.
Indeed, relative to the ample class $h_1+h_2$ on the surface
factor, the sum and difference directions have opposite
intersection signs.  Their Gram matrix is
$\operatorname{diag}(2,-2)$; changing the sign separately on the
primitive direction makes it positive.  The sum/difference change
of basis has determinant $-2$, so this observation is an isogeny
decomposition, not an integral direct-sum assertion at two.

The supplied program verifies smoothness of the displayed elliptic
curve, the intersection matrix, its determinant $-1$, and the
sum/difference diagonalization.  Its finite computations certify
exactly these algebraic assertions; the universal-property and
ampleness arguments are given above.

\paragraph{Known scope and a discarded integral-lattice route.}
The preceding example refutes positivity while possessing an
intersection isomorphism.  We next give a separate counterexample
to the universal isomorphism assertion itself.
For comparison, Achter--Casalaina-Martin--Vial prove a canonical
symmetric purely inseparable isogeny for geometrically rationally
chain connected threefolds, and an isomorphism for geometrically
stably rational threefolds; see \cite[Theorem 1, p.~6 and
Theorem 12.12, p.~54]{DGJ_ACMVDiagonal2025}.  Their Question 13.4,
p.~57, concerns positivity in the remaining nonliftable setting.
Those are stronger-hypothesis statements; they do not establish
the assertion of J.2 for arbitrary threefolds.

We tested a possible integral obstruction using
$E\times\operatorname{Km}(F\times G)$ over a finite field,
with distinct ordinary elliptic isogeny factors.  Its
Néron--Severi lattice can have discriminant sixteen, so the raw
cohomological restriction is not unimodular.  The attempted
deduction fails at identification of the integral Tate lattice
of $\operatorname{Ab}^2$: the assertion that the universal regular
homomorphism is an isomorphism on all prime-to-characteristic
torsion is a named \emph{standard assumption} in
\cite[Definition 6, p.~4; Proposition 3.8, p.~17]{DGJ_ACMVBloch},
not an unconditional theorem for all smooth threefolds in
positive characteristic.  The universal integral assertion in
the sweep's discussion of 44:8 must therefore not be used.
This route remains conditional; it is not a counterexample.


\subsection{A generic hyperelliptic Jacobian obstructs integral autoduality}
\paragraph{J2a (Refuted, with a complete counterexample).}
Let $a_0,\ldots,a_6$ be algebraically independent over $\mathbb F_5$,
put $K=\mathbb F_5(a_0,\ldots,a_6)^{\mathrm{perf}}$ and
$k=\overline K$, and let
\[
 C/K:\quad y^2=x^7+a_6x^6+\cdots+a_1x+a_0,
 \qquad A=\operatorname{Jac}(C),\qquad X=A_k.
\]
Here $C$ means the smooth projective model and $\infty$ is its
unique point above infinity.  Its genus is three, and $X$ is a
smooth projective threefold in characteristic five.  We prove
that there is no isomorphism from $\operatorname{Ab}^2(X)$ to its
dual induced by the integral intersection of families of cycles.
We use the precise required compatibility: for every family
$z\in\operatorname{CH}^2(T\times X)$, $z(0)=0$, on an abelian
parameter variety $T$, if $u_z:T\to\operatorname{Ab}^2(X)$ is its
regular homomorphism, the homomorphism defined by
$r_*(p^*z\cdot q^*z)$ must equal $u_z^\vee\rho u_z$.
This is exactly the construction in 44:9; no arbitrary rescaling
of intersection is allowed.

Write $B=\operatorname{Ab}^2(A/K)$.  This abelian variety exists,
and its base change to $k$ is the usual algebraic representative,
by \cite[Theorem 5.10, pp.~25--26; Theorem 6.1, pp.~26--27]{DGJ_ACMVFunctor}.
The compatibility defining $\rho$, if satisfied, determines it
uniquely: applying it to a miniversal family on $B$ gives
$[r]^\vee\rho[r]$ for some $r\ne0$, and $\operatorname{Hom}(B,B^\vee)$
is torsion-free.  Therefore every Galois conjugate of $\rho$ is the
same $\rho$.  It descends to $K$.  The perfectness of $K$ makes
$k/K$ separable, so the stated base-change and descent theorem
applies without an inseparable-field qualification.

\paragraph{The dimension of $B$ is at most three.}
Choose a miniversal family for $B$, and a prime $\ell\ne2,5$
not dividing its degree.  Its rational Tate realization is
injective by \cite[Proposition 8.2, p.~31]{DGJ_ACMVFunctor}.
After dualizing and applying Poincaré duality on the threefold,
this makes $H^1(B,\mathbb Q_\ell)$ a Galois quotient of
$H^3(A,\mathbb Q_\ell)(1)$.
The geometric Zariski monodromy on $V=H^1(A,\mathbb Q_\ell)$ is
$\operatorname{Sp}_6$.  Indeed, this is
\cite[Theorem 3.4 and Corollary 3.5, pp.~13--14]{DGJ_AchterPries2006}
for the hyperelliptic moduli space.  Here is the precise transfer to our coefficient space
$S=\operatorname{Spec}\mathbb F_5[a_0,\ldots,a_6,\Delta^{-1}]$.
Let $\mathcal H_3^w$ be the stack of smooth hyperelliptic
curves with one marked Weierstrass point.  The morphism
$\mathcal H_3^w\to\mathcal H_3$ forgetting the point is finite
étale of degree eight: its pullback to any family is that
family's reduced branch divisor.  Thus its geometric fundamental
group has finite-index image.  Finite-index subgroups of a
Zariski-dense subgroup of the connected group
$\operatorname{Sp}_6$ are still Zariski dense.

There is, more precisely, a presentation
$\mathcal H_3^w=[S/G]$ with connected smooth group
\[
 G=\{(b,a,c):b\in\mathbb G_a,\ a,c\in\mathbb G_m,
                    \ c^2=a^7\}
       \simeq\mathbb G_a\rtimes\mathbb G_m.
\]
The coordinate changes are $x\mapsto ax+b$, $y\mapsto cy$;
the monic leading coefficient imposes $c^2=a^7$.
The isomorphism on the multiplicative factor is
$s\mapsto(a,c)=(s^2,s^7)$, with inverse $s=c/a^3$.
Every pointed hyperelliptic curve has this form fppf locally:
the marked branch point trivializes the quotient conic locally,
we choose it as infinity, and we choose the anti-invariant
coordinate $y$ and normalize its leading coefficient.  Any
isomorphism of two such presentations has exactly the displayed
form, since functions with a pole of order at most two at the
marked point are spanned by $1,x$, and the anti-invariant
functions with pole order at most seven are spanned by $y$.
This proves the claimed presentation, also for families by the
corresponding locally free sheaves and base change.
The atlas $S\to[S/G]$ has geometrically connected fibres $G$,
so it surjects on geometric fundamental groups: a connected
finite étale cover of the stack stays connected after this
pullback, which is the finite-cover criterion for surjectivity.
Consequently the geometric Zariski monodromy on $S$ is
$\operatorname{Sp}_6$.  Removing a proper closed subset of
normal $S$ gives a surjection on fundamental groups, and
purely inseparable extension to $K$ leaves the étale Galois
representations unchanged.

The symplectic exterior-cube decomposition is
\[
 \bigwedge^3V(1)=V\oplus P,
 \qquad \dim V=6,\quad \dim P=14,
\]
where the displayed $V$ denotes its Lefschetz copy and $P$ is
primitive.  Both summands are absolutely irreducible for the
geometric symplectic group and are preserved by the arithmetic
Galois action.  A direct verification of the only less familiar
assertion is as follows.  With symplectic basis $e_i,f_i$, $P$
is the kernel of symplectic contraction.  Its fourteen distinct
torus weights consist of the eight $\pm\epsilon_1\pm\epsilon_2
\pm\epsilon_3$ and the six $\pm\epsilon_i$.  The first weight
vectors are wedges choosing one of $e_i,f_i$ for each $i$; the
second are $v_i\wedge(e_j\wedge f_j-e_h\wedge f_h)$ for
$\{i,j,h\}=\{1,2,3\}$ and $v_i=e_i$ or $f_i$.
Symplectic permutations and interchanges generate all eight
vectors of the first kind from any one.  For an explicit root operation, let $N(e_2)=e_3$,
$N(f_3)=-f_2$, and let $N$ vanish on the remaining basis
vectors.  It belongs to the symplectic Lie algebra and sends
$e_1\wedge e_2\wedge f_3$ to
$-e_1\wedge(e_2\wedge f_2-e_3\wedge f_3)$.
The opposite root $N'(e_3)=e_2$, $N'(f_2)=-f_3$ sends the
latter second-kind vector to $-2e_1\wedge e_2\wedge f_3$.
Permuting and interchanging symplectic pairs gives all six
second-kind vectors and the converse transitions.  Thus a nonzero
subrepresentation, which contains a torus weight vector, contains
all fourteen vectors.  This proves irreducibility over
$\mathbb Q_\ell$, where $2\ne0$.
Consequently the quotient $H^1(B)$ can have only the constituents
$V$ and $P$, each at most once.

The constituent $P$ cannot occur.  We give the spreading
step explicitly.  Write
$K_n=\mathbb F_5(a_0^{1/5^n},\ldots,a_6^{1/5^n})$, so
$K=\bigcup_nK_n$.  The abelian variety $B$, its group laws,
and the finitely many components and coefficients of a
miniversal cycle $\Gamma\in\operatorname{CH}^2(B\times A)$
are of finite presentation; hence they descend to some $K_n$.
The model of $A$ is the Frobenius pullback of its model over
$S$.  After removing a proper closed subset of this Frobenius
pullback $S_n$, the descended $B$ extends to an abelian scheme
$\mathcal B\to U$, and the closures of the cycle components
give a cycle on $\mathcal B\times_U\mathcal A$.  To obtain
the abelian scheme, extend projective equations and the group
laws and their inverse, shrink so that the group identities
hold, and shrink again to the smooth locus with geometrically
connected fibres.  Properness is retained by the projective
model.  No extension of the universal regular property to every
fibre is being assumed.

The generic Tate injection of Proposition~8.2 is equivariant
for $\operatorname{Gal}(K_n)$: purely inseparable base change
identifies this group and its étale cohomology with those over
$K$.  Both its source and target are the generic fibres of
lisse $\mathbb Q_\ell$-sheaves on $U$, and therefore factor
through $\pi_1(U)$.  Since
$\operatorname{Gal}(K_n)\to\pi_1(U)$ is surjective, their
equivariant map is a map of those lisse sheaves.  Its rank is
locally constant.  Dualizing gives a surjection on every fibre
of $U$, just as at the generic point.  Its kernel is generically
one of the sums of the two intrinsic Lefschetz summands;
equality of morphisms of lisse sheaves at a geometric generic
point is equality everywhere, so its fibrewise quotient has the
same summands.  The Frobenius pullback and shrinking preserve
connected geometric monodromy by the preceding fundamental-group
argument.  The ordinary locus of the Jacobian is
nonempty.  For an explicit check, the smooth curve
$y^2=x^7-x$ over $\mathbb F_5$ has Cartier--Manin matrix
\[
 \begin{pmatrix}0&0&1\\0&3&0\\1&0&0\end{pmatrix},
 \qquad\det=2\pmod5.
\]
This follows directly from applying Cartier to
$x^{j-1}dx/y=y^{-5}x^{j-1}(x^7-x)^2dx$:
the $(i,j)$ entry is the coefficient of $x^{5i-j}$ in
$(x^7-x)^2$.  Thus Cartier is invertible, so the Jacobian is
ordinary.  Nonvanishing of the same determinant defines a
nonempty open ordinary locus in the coefficient space.  It
meets the open set where the preceding objects have spread;
choose a closed point there, over a finite field $\mathbb F_q$.

For clarity, the passage from the Cartier calculation to
Frobenius slopes can be made entirely in terms of characteristic
polynomials.  The Manin congruence in
\cite[\S2.5, p.~7; Cartier formula in \S3.1, pp.~7--8]{DGJ_AchterHowe2018}
says that the Frobenius polynomial modulo five is $T^3$ times
the characteristic polynomial of an invertible iterate of the
Cartier operator (equivalently its adjoint Frobenius operator).
That second polynomial has nonzero constant coefficient.
Hensel factorization therefore gives exactly three unit roots.
The principal polarization pairs each root $\alpha$ with
$q/\alpha$; see also \cite[Lemma 19.2, p.~41]{DGJ_MilneAV}.
Hence the remaining three roots have normalized valuation one.
Thus the six Frobenius eigenvalues on $H^1(A)$ have three
$q$-adic slopes zero and three slopes one.  Their triple
products on $H^3(A)(1)$ include one slope $-1$ and one slope $2$.
The Lefschetz summand has only slopes zero and one, so both
extreme slopes belong to $P$.  The surjection on $\ell$-adic fibres gives equality of
Frobenius characteristic polynomials for $H^1(B)$ and the
indicated summand or sum of summands.  These polynomials have
rational coefficients: the primitive polynomial is the quotient
of the exterior-cube polynomial by the Lefschetz polynomial.
Equality in $\mathbb Q_\ell[T]$ is therefore equality in
$\mathbb Q[T]$.  Embed their algebraic roots in
$\overline{\mathbb Q}_5$ and normalize $v(q)=1$.
If $P$ occurred in the quotient $H^1(B)$, its negative-valuation
root would then be a Frobenius eigenvalue on $H^1(B)$.  This
contradicts the algebraic integrality of Frobenius eigenvalues
on abelian varieties, as stated in
\cite[Theorem 19.1 and its proof, pp.~40--41]{DGJ_MilneAV}.
This argument compares characteristic polynomials; it makes no
assertion that the correspondence admits a particular crystalline
realization.  Hence only the
six-dimensional constituent can occur, and $\dim B\leq3$.
The finite computation records the full primitive slope
multiplicities $1,6,6,1$ in slopes $-1,0,1,2$; the exclusion uses
only the first of these.

\paragraph{The translation family has intersection $\pm2\lambda_A$.}
Embed $C$ by $Q\mapsto[Q-\infty]$ and let
\[
 z_t=[C+t]-[C],\qquad t\in A.
\]
Equivalently $z=m_-^*[C]-\operatorname{pr}_X^*[C]$ on
$A\times A$, where $m_-(t,x)=x-t$.  This is a specified family
of algebraically trivial codimension-two cycles with $z_0=0$.
It induces a $K$-homomorphism $u:A\to B$.
Let $\theta$ denote the cohomology class of the principal theta
polarization.  The Poincaré formula is $[C]=\theta^2/2$ in
cohomology.  It follows as well by characterizing $[C]$ by
$\int_A[C]\alpha\beta=\int_Ca^*\alpha\,a^*\beta$ and the
symplectic pairing defining the principal Jacobian polarization.
With symplectic basis $e_i,f_i$ and
$\theta=\sum e_i\wedge f_i$, the part of the class of $z$ in
$H^1(A_{\rm par})\otimes H^3(A_{\rm space})$ is obtained by
expanding
\[
 \frac12\left(\sum_{i=1}^3
 (e_{i,x}-e_{i,t})\wedge(f_{i,x}-f_{i,t})\right)^2
       -\frac12\left(\sum e_{i,x}\wedge f_{i,x}\right)^2.
\]
All terms of parameter degree zero cancel.  Therefore in
$r_*(p^*z\cdot q^*z)$ only the product of the two parameter-degree-one
parts contributes.  Wedge multiplication and integration of
$e_{1,x}f_{1,x}e_{2,x}f_{2,x}e_{3,x}f_{3,x}$ give
\[
 -2\sum_{i=1}^3
 (e_{i,t}\wedge f_{i,s}-f_{i,t}\wedge e_{i,s}).
\]
This is $-2$ times the normalized Poincaré pairing of the
principal polarization, with the sign reversed if the opposite
Poincaré convention is chosen.  Faithfulness of the Tate
realization on homomorphisms of abelian varieties
\cite[Theorem 12.5, pp.~21--22]{DGJ_MilneAV} gives the
actual homomorphism $\pm2\lambda_A$, not merely its isogeny class.
The wedge calculation is also saved as an exact integer
calculation in \texttt{check\_generic\_jacobian.py}.

\paragraph{The impossible degree.}
If an intersection-induced isomorphism $\rho:B_k\to B_k^\vee$
existed, its defining compatibility and the preceding calculation
would give
\[
 u^\vee\rho u=\pm2\lambda_A.
\]
The right side is an isogeny.  Thus $u$ has finite kernel;
since $\dim B\leq3=\dim A$, it is an isogeny onto $B$.
Taking degrees yields
\[
 (\deg u)^2=\deg(2\lambda_A)=2^6,
 \qquad\deg u=8.
\]
We show that the $K$-abelian variety $A$ admits no isogeny of
degree eight.  The splitting field of its generic septic has
Galois group $S_7$: equivalently, start with seven independent
roots and take the invariant field of their permutations.
Passing to the perfect closure is purely inseparable and leaves
this group unchanged.  If $r_1,\ldots,r_7$ are the roots and $W_1,\ldots,W_7$ are the finite
Weierstrass points, the elements $[W_i-\infty]$ generate $A[2](k)$,
with their sum as the only relation.  To check the last assertion,
a putative proper-subset relation would make
$\prod_{i\in I}(x-r_i)$ a square in $k(C)$.  In the quadratic
extension $k(C)/k(x)$, a rational function whose square lies in
$k(x)$ is in $k(x)$ or in $y\,k(x)$.  The odd valuations at the
chosen roots exclude the first case unless $I$ is empty and
exclude the second unless $I$ is the full set.  There are
therefore six independent classes, as required by the rank of
$A[2]$.

It follows that $A[2](k)$ is the six-dimensional $S_7$-module
$\mathbb F_2^7/\langle(1,\ldots,1)\rangle$, equivalently the
subspace of vectors with even coordinate sum.  This module is
irreducible.  Indeed any nonzero even-sum vector has both a zero
and a one coordinate; subtracting its image under the
transposition of those coordinates yields a vector supported
on exactly two positions.  The $S_7$-orbit of that vector spans
the whole even-sum subspace.  Every nonzero invariant subspace
is consequently the whole module.

The kernel of a degree-eight isogeny in characteristic five is
an étale group of order eight.  Its subgroup killed by two is a
nonzero Galois-stable subgroup of $A[2](k)$, and has at most
eight elements.  Irreducibility would force it to be all of
$A[2](k)$, which has sixty-four elements.  This contradiction
proves the claimed impossibility of an integral
intersection-induced isomorphism.

This argument uses only the rational miniversal Tate embedding,
not the unproved integral torsion standard assumption.  It
rules out the raw integral autoduality demanded in J.2.  It does
not rule out a rational pairing, an isogeny, a polarization after
an appropriate Lefschetz normalization, or an unrelated abstract
isomorphism of $B$ with its dual.


\subsection{An explicit counterexample to numerical equals tau equivalence}
\paragraph{J1-Th7 (Refuted).}
The manuscript statement at 44:6 is
``L'équivalence numérique $=$ $\tau$-équivalence''.
Its valid general direction is that $\tau$-equivalence implies
numerical equivalence: intersection numbers are constant in a
connected algebraic family, and dividing by a nonzero integer
does not change vanishing of integer intersection numbers.
The reverse implication fails already for codimension two on
an abelian threefold over $\overline{\mathbb Q}$.

Here are all data of one counterexample.  Put
\[
 k=\overline{\mathbb Q},\qquad
 C=V(x^4+y^4-z^4)\subset\mathbb P^2_k,\qquad
 P=(1:0:1),\qquad X=\operatorname{Jac}(C).
\]
The three partial derivatives have no simultaneous projective
zero, so $C$ is a smooth plane quartic of genus three.  In
particular $X$ is a smooth projective threefold.  Let
$a_P:C\hookrightarrow X$ be $Q\mapsto[Q-P]$, and define
\[
 Z=(a_P)_*[C]-[-1]_* (a_P)_*[C]\in\operatorname{CH}^2(X).
\]
The action of $[-1]$ on $H^4(X,\mathbb Q_\ell(2))$ is $+1$,
since the cohomology algebra of an abelian variety is the
exterior algebra of $H^1$ and $[-1]$ acts by $-1$ on $H^1$.
Thus the class of $Z$ is zero.  Compatibility of cycle classes
with products and degree gives $\deg(Z\cdot D)=0$ for every
divisor $D$ on $X$.  This is precisely numerical triviality
of this codimension-two cycle.

Bloch's \emph{Algebraic cycles and values of L-functions},
introduction p.~95 and Theorem~(4.1), p.~104
\cite{DGJ_Bloch1984}, proves that the class of this very
Fermat-quartic cycle has infinite order modulo algebraic
equivalence.  His definition of that quotient is on p.~94,
and p.~95 states that its class is independent of the chosen
Abel embedding.  The proof on p.~104 works over an algebraic
closure of the number field, so it applies to the field specified
above; no passage to a generic curve or unspecified very general
parameter is involved.  Consequently there is no nonzero integer
$N$ for which $NZ$ is algebraically trivial.  Equivalently,
$Z$ is not $\tau$-equivalent to zero.  The use of the published
theorem is exactly the non-algebraicity step; the data,
smoothness, homological triviality, and numerical consequence
have been supplied explicitly.

For completeness the same example proves strictness in every
intermediate codimension $2\leq i\leq n-1$ in every dimension
$n\geq3$.  Put $m=n-3$, $r=i-2$, $X_n=X\times\mathbb P^m$,
and $Z_i=\operatorname{pr}_X^*Z\cdot h^r$.  It is numerically
trivial by the projective bundle formula.  If a nonzero multiple
of $Z_i$ were algebraically trivial, intersecting it with
$h^{m-r}$ and pushing forward to $X$ would show that the same
multiple of $Z$ were algebraically trivial, a contradiction.
The operations used preserve algebraic equivalence.  Hence the
correct unrestricted comparison is a strict implication in
these intermediate degrees, not equality.

The saved calculation \texttt{check\_fermat.py} verifies the
singular ideal is the unit ideal in each of the three projective
charts, verifies the specified point, and records genus three.
It does not purport to compute Bloch's infinite-order theorem.
The original scanned source pages were rendered and visually
read, including pp.~94--95 and 104; the complete volume PDF
has an offset of seven pages from its printed numbering.


\subsection{SA19 H6: what the planning leaves establish}
\paragraph{H6a (Textual and variance correction, proved).}
The morphism from an open curve is explicitly present:
43:27 boxes $\varphi_S:X-S\to\mathscr P^1_{S,X/Y}$.
At 43:28 its universal property is called the theorem of
Rosenlicht.  Corollary~3 at 43:30 states an ``hom. canonique
\emph{injectif}''
\[
 \operatorname{Ext}^1(\mathscr P_{S,X/Y},G)
        \longrightarrow H^1(X-S,G).
\]
The next sentence says only ``Il \emph{semble} que c'est
\emph{[bijectif|surjectif]}'' when $S\to Y$ is surjective.
The supplied transcription leaves the bracketed word uncertain;
either reading describes the conjectured stronger direction.
43:31 proposes reductions, and 43:34 writes ``et sans doute''
for the higher-degree comparison.  These qualifications must be
retained in a documentary assessment.

The sweep's stated inference to injectivity on abelian
fundamental groups has its variance reversed.  When the two
sides are interpreted as finite-constant torsors, the map is
pullback of characters:
\[
 \operatorname{Hom}_{\rm cont}(\pi_1(J)^{\rm ab},M)
 \longrightarrow
 \operatorname{Hom}_{\rm cont}(\pi_1(U)^{\rm ab},M).
\]
Its injectivity for every finite abelian $M$ detects
\emph{surjectivity} of $\pi_1(U)^{\rm ab}\to\pi_1(J)^{\rm ab}$.
Indeed a proper closed image in a profinite abelian group is
detected by a nonzero character of a finite quotient, and
conversely a surjection gives injective pullback.  It says
nothing by itself about the kernel.  The quotient
$\mathbb Z/4\to\mathbb Z/2$ has injective pullback of characters
for every finite coefficient group and a nonzero kernel.
If the stronger bijectivity were separately proved, it would
yield an isomorphism of the represented profinite groups.
Neither conclusion is a point-to-section injectivity theorem.

\paragraph{H6b (Complete scoped documentary audit).}
The exact lead extends its negative finding to the full range
43:1--208. The audit now covers all 195 supplied nonblank
transcription files: DGJ read 184 complete files, and FHI
explicitly certified a complete rereading of the eleven files
43:56--66. No readable rational-point-to-section injectivity
argument was found. The saved per-file inventory records source
paths, SHA-256 hashes, topic groups, readers, and uncertainty
counts in \texttt{work/dgj/outputs/c43\_audit.tsv}; the complete
scope discussion is \texttt{work/dgj/audit\_SA19H6.md}.
The earlier 71 keyword matches served only to plan the reading.

The near matches have different subjects. At 43:77 a rational
section of a torsor extends; at 43:103 the reused printed text
concerns cross-sections of Stiefel manifolds. At 43:145 the
injectivity is a localization statement in $H^1$. At
43:167--170 a normalized torsor section proves injectivity from
group extensions to primitive torsor classes. At 43:204--205,
sections of a dual abelian scheme define the Tate pairing;
43:207 uses the Albanese map to define a zero-cycle kernel and
a pairing. None is the sought rational-point-to-section map.

This finishes the \emph{readable supplied-transcription} audit.
It does not prove absence in original handwriting, unresolved
alternate readings, illegible strings, unsupplied material, or
other archival items. The corpus retains 2,620 uncertainty and
48 illegibility macros. Thirteen scan labels have no separate
file: 14, 39, 44, 83, 88, 99, 111, 127, 157, 166, 175, 179,
and 194; their contents are not inferred. Thus the scoped
negative finding is complete, while a stronger historical
nonexistence assertion remains unestablished.

% BEGIN J5 completed quadric counterexample
\subsection*{J5: the height formula fails for a split quadric fourfold}
\textbf{Status: Refuted.} Root independently audited the complete proof and both supplied transcriptions of 44:4.
The failure below concerns the literal statement of Theorem5, 44:4.
Its insertion reads ``section hyperplane générique d'un pinceau dont
[les] fibres sont géom. intègres'', and it concludes
$\langle\zeta,\zeta'\rangle=(-1)^mZ\mathbin{.}Z'$.
The example satisfies this statement,
including the inserted requirement that all pencil fibres be
geometrically integral. It does not depend on identifying the
manuscript's equivalence relation or constructing any nonzero height:
both restricted cycles are literally zero. The original middle
case $n=2m=4$, $i=j=m=2$ also makes its displayed sign unambiguous.

Let $k=\overline{\mathbb F}_3$, and let
\[
 X=V(x_0x_3+x_1x_4+x_2x_5)\subset\mathbb P^5_k.
\]
The six partial derivatives are $x_3,x_4,x_5,x_0,x_1,x_2$,
so $X$ is a smooth projective fourfold. Consider the pencil
\[
 H_t=V(x_2-tx_5),\quad t\in\mathbb A^1,
 \qquad H_\infty=V(x_5).
\]
Its base locus in $X$ is the smooth quadric surface
$B=V(x_2,x_5,x_0x_3+x_1x_4)\subset\mathbb P^5$.
Blowing up $B$ resolves the pencil and gives a smooth projective
total space over $\mathbb P^1$; the fibres remain $X\cap H_t$,
because $B$ is a Cartier divisor in each member. For finite $t$,
the fibre has equation
\[
                 x_0x_3+x_1x_4+t x_5^2=0
                 \quad\text{in }\mathbb P^4.
\]
It is smooth for $t\ne0$ and is a rank-four quadric cone for $t=0$.
The fibre at infinity is another rank-four quadric cone, with
equation $x_0x_3+x_1x_4=0$ and extra coordinate $x_2$.
All fibres are geometrically integral: a quadratic form which is
a product of two linear forms has rank at most two, whereas these
have rank at least four, over every algebraic extension.
In particular this is a Lefschetz pencil with smooth base and two
ordinary quadratic singular fibres, rather than a pencil whose
total family is degenerate.

Take the four planes
\[
 \alpha=V(x_3,x_4,x_5)\subset X,
 \qquad
 \beta=V(x_2,x_3,x_4)\subset X,
\]
\[
 \alpha'=V(x_0,x_1,x_5)\subset X,
 \qquad
 \beta'=V(x_0,x_1,x_2)\subset X.
\]
Put $Z=[\alpha]-[\beta]$, $Z'=[\alpha']-[\beta']$.
If $K=k(t)$ and $Y=X\cap H_t$ is the generic
member, then the proper scheme-theoretic intersections are
\[
 \alpha\cap Y=\beta\cap Y
   =L:=V(x_2,x_3,x_4,x_5)\simeq\mathbb P^1_K,
\]
and
\[
 \alpha'\cap Y=\beta'\cap Y
   =L':=V(x_0,x_1,x_2,x_5)\simeq\mathbb P^1_K,
\]
all with multiplicity one. Indeed $H_t$ restricts on $\alpha$
to $x_2=0$, and on $\beta$ to $-t x_5=0$, where $t$ is a unit
in $K$; the same statements hold for $\alpha',\beta'$.
Therefore
\[
                  \zeta=Z\cdot Y=0,
                  \qquad \zeta'=Z'\cdot Y=0             \tag{J5.1}
\]
already in the group of cycles, not only modulo an equivalence.
Any Neron--Tate bilinear pairing on their associated points is zero.

Here is a direct calculation using only proper transverse
intersections. The crossed pairs $\alpha\cap\beta'$ and
$\beta\cap\alpha'$ are empty. The other two intersections are
\[
 \alpha\cap\alpha'=\{[e_2]\},\qquad
 \beta\cap\beta'=\{[e_5]\}.
\]
At $[e_2]$, the chart $x_2=1$ writes the smooth fourfold with local
coordinates $x_0,x_1,x_3,x_4$ and
$x_5=-x_0x_3-x_1x_4$. The ideals of the two planes are
$(x_3,x_4)$ and $(x_0,x_1)$. They therefore intersect transversely
with local multiplicity one. At $[e_5]$, the chart $x_5=1$
gives exactly the same coordinate calculation, with
$x_2=-x_0x_3-x_1x_4$. Consequently
\[
                  Z\cdot Z'=1+1-0-0=2.                 \tag{J5.2}
\]
The integer two in (J5.2) is the degree of a zero-cycle; it is not
to be reduced modulo the characteristic.

Equations (J5.1)--(J5.2) contradict exactly the proposed formula
\[
             \langle\zeta,\zeta'\rangle
                       =(-1)^m Z\cdot Z',
                  \qquad m=2,
\]
by giving $0=2$. Changing one overall sign cannot fix it.
Both cycles have algebraically trivial generic restrictions, and
every member of the pencil is geometrically integral, as required
by the revised ink. Thus the literal higher-codimension Theorem5
is false even in this smallest balanced case beyond surfaces.

\paragraph{What the integral-fibre condition misses.}
The obstruction is exact: generic restriction does not determine
intersection numbers of higher-codimension extensions. The class
$Z$ is in the kernel of generic restriction but pairs nontrivially with $Z'$,
so this intersection pairing cannot descend to that quotient.
In this displayed example each plane meets the pencil base in a line,
so those intersections have excess dimension. The literal statement
does not forbid this. An additional condition on extensions or their
incidence with the base would be a stronger assertion requiring a
separate analysis. After blowing up the pencil base, the strict transform of $\alpha$
lies in the fibre at infinity and that of $\beta$ in the fibre at
zero; their common exceptional correction cancels in the difference.
Geometric integrality eliminates independent vertical divisor
components, but it does not eliminate these vertical codimension-two
classes through the singularities of the quadric fibres.

The following precise elementary replacement records the necessary
correction. For any rational vector space $V$ with a bilinear form
$b$ and a restriction map $r:V\to W$, let $K=\ker r$.
The form factors through $r(V)$ if and only if
$b(K,V)=b(V,K)=0$. Necessity follows by pairing a zero image with
every image; sufficiency follows by changing either lift by an
element of $K$. For a symmetric form, its restriction always
descends to
\[
                K^\perp/(K\cap K^\perp).
\]
Apply this to rational cycle classes on a smooth projective model,
their intersection form, and restriction to the generic fibre.
A higher-codimension height comparison needs an admissibility or
vertical-correction condition that supplies this orthogonality;
geometric integrality of the fibres alone does not supply it.
This counterexample does not refute a comparison with such
additional conditions, nor the classical divisor comparison for
surfaces with integral fibres.

\paragraph{Finite check.}
The saved script \texttt{check\_quadric\_height.py} checks the
gradient, the pencil ranks, all four literal restriction ideals,
the two transverse point intersections and the two empty intersections,
and the integer intersection value. An additional normal-bundle and
moving-plane calculation independently checks $Z^2=2$ as well.
The written proof supplies the geometric and height conclusions.

% END J5 completed quadric counterexample

% BEGIN J6L independently audited conditional theorem
\subsection{A precise conditional reduction for the Lefschetz assertions}
\paragraph{J6L (Conditional theorem, proved).}
Let $k$ be algebraically closed, of any characteristic. Let
$j:Y\hookrightarrow X$ be a smooth hyperplane section of a smooth
projective $n$-fold, and use the incidence abelian varieties of J1i.
Fix $\ell\ne\operatorname{char}k$ and put $r=2i-1$. For a smooth
projective $d$-fold $W$, write $L_W$ for cup product by the chosen
hyperplane class. The convention for the Lefschetz lowering
operator is
\[
 \Lambda_W(L_W^a v)=
 \begin{cases}L_W^{a-1}v,&a>0,\\0,&a=0,\end{cases}
 \qquad v\in H^b_{\rm prim}(W,\mathbb Q_\ell).
\]
Thus $\Lambda_W$ lowers degree by two and Tate twist by one; it
has no extra $\mathfrak{sl}_2$ coefficient. Let $B(W)$ mean that
this operator is induced by a rational algebraic correspondence
in $\operatorname{CH}^{d-1}(W\times W)_{\mathbb Q}$.
With this explicit convention the following implications hold:
\begin{enumerate}
\item If $2i\leq n-1$ and $B(X)$ holds, then
$j^*:P^i_{\rm inc}(X)\to P^i_{\rm inc}(Y)$ is an isogeny.
\item If $2i\geq n+1$ and $B(Y)$ holds, then the same homomorphism
is surjective.
\end{enumerate}
Only the relevant degree of $\Lambda_X$, respectively
$\Lambda_Y$, is needed. The cohomological inputs are Deligne's hard Lefschetz theorem
\cite[Theorem 4.1.1, p.~217]{DGJ_Deligne1980} and weak
Lefschetz \cite[(4.1.6), pp.~218--219]{DGJ_Deligne1980}.
The coefficient-free lowering convention is also spelled out in
\cite[Section 4.1, pp.~6--7]{DGJ_Peters2026}.
No integral inverse, equality of
incidence and universal regular equivalence, or middle-degree
fixed-part assertion is part of this conclusion.

\emph{Cohomological detection lemma.}
If $U,V$ are smooth projective, $d=\dim U$, and
$\Gamma\in\operatorname{CH}^{d+b-a}(U\times V)_{\mathbb Q}$
acts as zero on
\[
 H^{2a-1}(U,\mathbb Q_\ell(a))
       \longrightarrow H^{2b-1}(V,\mathbb Q_\ell(b)),
\]
then the induced rational homomorphism
\[
 P^a_{\rm inc}(U)\longrightarrow P^b_{\rm inc}(V)
\]
is zero. Here a rational homomorphism is an element of
$\operatorname{Hom}_{k\text{-av}}(-,-)\otimes\mathbb Q$.
To prove the lemma, set $A=P^a_{\rm inc}(U)$, $e=\dim A$.
Choose a miniversal cycle
$u\in\operatorname{CH}^a(A\times U)$ inducing $[m]$ on $A$,
where $m>0$; see \cite[Lemma 2.3, pp.~104--105]{DGJ_Saito1979}.
Choose a Picard realization of the universal map for $V$, namely
a correspondence $w$ and a closed immersion
$P^b_{\rm inc}(V)\hookrightarrow\operatorname{Pic}^{\rm var}(T)$,
as in \cite[(3.5), p.~110]{DGJ_Saito1979}. The composite
\[
 D=w\circ\Gamma\circ u\in
               \operatorname{CH}^1(A\times T)_{\mathbb Q}
\]
induces the composite of $[m]$, the proposed homomorphism, and
this immersion. Its action on
$H^{2e-1}(A,\mathbb Q_\ell(e))\to
H^1(T,\mathbb Q_\ell(1))$ is zero by the hypothesis.
For a divisor correspondence on $A\times T$, this action is the
Poincaré-dual form of its homomorphism on rational Tate modules.
Consequently the displayed composite homomorphism is zero:
a nonzero homomorphism of abelian varieties has a positive
dimensional abelian-variety image and hence has nonzero rational
Tate-module image. The closed immersion and the invertibility
of $[m]$ in the isogeny category give the lemma. Equivalently one
may clear all denominators first; torsion-free Hom groups give
the same conclusion. The case $A=0$ is immediate.

We now prove the two assertions. The cohomology conventions give
\[
 j_*j^*=L_X,\qquad j^*j_*=L_Y,
 \qquad
 j_*:H^q(Y,\mathbb Q_\ell)
         \longrightarrow H^{q+2}(X,\mathbb Q_\ell(1)).
\]
In the first range, $r\leq n-2$, weak Lefschetz makes $j^*$ an
isomorphism on $H^r$. Every Lefschetz string in $H^r(X)$ has
another step above it, so $\Lambda_X L_X=1$ there. The rational
correspondence
\[
 \beta=\Lambda_X\circ j_*\in
           \operatorname{CH}^{n-1}(Y\times X)_{\mathbb Q}
\]
therefore induces the inverse of $j^*$ on $H^r$:
$\beta j^*=\Lambda_XL_X=1$, and $j^*\beta=1$ follows since
$j^*$ is an isomorphism. The detection lemma applied to the two
composites minus the diagonals proves that $\beta$ is a two-sided
inverse on incidence abelian varieties in the isogeny category.
After clearing denominators, each composite is multiplication
by the same positive integer. Thus the actual restriction
homomorphism is an isogeny.

In the second range, $r\geq n>\dim Y$, every Lefschetz string in
$H^r(Y)$ has a step below it. Hence $L_Y\Lambda_Y=1$ there.
The rational correspondence
\[
 \beta=j_*\circ\Lambda_Y\in
           \operatorname{CH}^{n-1}(Y\times X)_{\mathbb Q}
\]
satisfies $j^*\beta=L_Y\Lambda_Y=1$ on $H^r(Y)$. Detection
therefore gives a right inverse in the isogeny category.
Clearing denominators makes the composite multiplication by a
positive integer on $P^i_{\rm inc}(Y)$, so the actual
restriction is surjective. If $i>\dim Y$, the target is zero
and surjectivity is automatic. This proves J6L.

The proof specifies exactly what cohomological weak Lefschetz
alone lacks: algebraic correspondences for the relevant lowering
operators. It also explains why the stable-range conclusion here
is an isogeny and does not settle the alternate transcription
``isomorphisme''. The middle fixed-factor clause and removal of
the stated algebraicity hypotheses remain active gates.

For comparison, in characteristic zero there is a further precise
reduction relating incidence equivalence to the algebraic part of
the Abel--Jacobi map. Peters' Corollary~7.8, p.~16, states that
if all primitive Lefschetz components of the relevant algebraic
odd cohomology remain algebraic, then these two quotients agree
up to isogeny \cite{DGJ_Peters2026}. Standard conjecture $B(X)$
implies that primitive-component condition; the codimension-two
condition is known, as summarized in his Section~7, pp.~15--16.
This result concerns complex Abel--Jacobi and incidence quotients.
It does not identify the universal regular quotient in all
codimensions and does not remove the positive-characteristic
algebraicity gate in J6L.

% END J6L

% BEGIN J6P: smooth-axis trace theorem and top-target surjectivity
\subsection{Codimension-two pencils and the trace image}
\paragraph{J6P (Proved with the stated smooth-axis hypothesis).}
The middle assertion at 44:4 reads ``le facteur fixe de $P^i(Y)$
est identique à l'image de $P^i(X)$''. Here $\operatorname{Ab}^2$ means the universal regular algebraic
representative, not an unidentified manuscript quotient.
Let $k$ be algebraically closed, let $X/k$ be smooth projective of
dimension $n\geq4$, and let a hyperplane pencil have a smooth
codimension-two base locus $B$. Assume its generic member $Y/K$,
$K=k(t)$, is smooth. Write $j:Y\hookrightarrow X_K$ and
$A=\operatorname{Ab}^2_{Y/K}$. Then
\begin{equation}\label{eq:J6P-trace}
 \operatorname{im}\bigl(j^*:\operatorname{Ab}^2_{X/k,K}\longrightarrow A\bigr)
 =\operatorname{im}\bigl(\tau:T_K\longrightarrow A\bigr),
 \qquad T=\operatorname{Tr}_{K/k}(A),
\end{equation}
as abelian subvarieties of $A$. If $n\geq5$, both images are $A$,
so $j^*$ is surjective. No standard conjecture or torsion
injectivity assumption is used for these two conclusions.
The smoothness of the pencil axis is an additional hypothesis:
it is not inferred from the manuscript's integral-fibre condition.

The representative and its separable base-change compatibility
are supplied by \cite[Theorems 6.1 and 5.10, pp.~25--27]{DGJ_ACMVFunctor}.
The trace exists by \cite[Definition 6.1 and Theorem 6.2, p.~20]{DGJ_Conrad2006}.
We use its image, since its kernel can be infinitesimal in positive
characteristic; see \cite[Theorem 6.12, pp.~22--23]{DGJ_Conrad2006}.

\emph{Step 1: spreading a trace family.}
Set $\widetilde X=\operatorname{Bl}_B X$, with blowdown $p$,
exceptional divisor $e:E\hookrightarrow\widetilde X$, and
$q:E\to B$. The pencil gives $f:\widetilde X\to\mathbb P^1$.
Its generic fibre is $Y$: blowing up a member along its Cartier
divisor $B$ changes nothing. Both $\widetilde X$ and $B$ are
smooth projective over $k$.

Choose a miniversal cycle
$\Gamma\in\operatorname{CH}^2(A\times_KY)$ whose induced map
under the universal regular homomorphism $\Phi_Y$ is $[r]_A$,
where $r>0$. This exists over $K$ itself, including when $K$ is
imperfect, by \cite[Lemma 4.7, pp.~18--19]{DGJ_ACMVFunctor}.
Pull it back along $\tau\times1$ to
$u\in\operatorname{CH}^2(T_K\times_KY)$. Normalize at the zero
of $T$ if necessary. Then
\[
 \Phi_Y(u_t-u_0)=r\tau(t).
\]
Take closures of cycle representatives to obtain
$Z\in\operatorname{CH}^2(T\times_k\widetilde X)$ restricting
to $u$ over the generic point of $\mathbb P^1$. Refined fibres
$Z_t-Z_0$ are algebraically trivial and give a regular
homomorphism $v:T\to\operatorname{Ab}^2_{\widetilde X/k}$.
The generic restriction $\delta$ therefore satisfies
\begin{equation}\label{eq:J6P-spread}
 \delta\circ v_K=r\tau,
 \qquad
 \delta:\operatorname{Ab}^2_{\widetilde X/k,K}\longrightarrow A.
\end{equation}
This equality holds as morphisms: it is the equality of regular
maps induced by the generic restriction of the same Chow class.

\emph{Step 2: the exceptional summand.}
The codimension-two blowup formula is
\[
 \operatorname{CH}^2(\widetilde X)
 =p^*\operatorname{CH}^2(X)\oplus e_*q^*\operatorname{CH}^1(B).
\]
Its inverse is $(p_*,-q_*e^*)$. Indeed,
$p_*p^*=1$, $p_*e_*q^*=0$, $q_*e^*p^*=0$, and
$q_*e^*e_*q^*=-1$, the last identity following from
$N_{E/\widetilde X}=\mathcal O_E(-1)$ and
$q_*c_1(\mathcal O_E(1))=1$. The general formula is opened in
\cite[Section 2.3, Theorem 2.8(3)]{DGJ_Jiang2023}, which states
Fulton's blowup formula. All these maps preserve algebraic
equivalence and commute with the refined fibres in a smooth
parameter space. Consequently the decomposition is compatible
with regular homomorphisms and gives an actual isomorphism
\begin{equation}\label{eq:J6P-blowup}
 \operatorname{Ab}^2_{\widetilde X/k}
 \simeq\operatorname{Ab}^2_{X/k}\times\operatorname{Pic}^{\rm var}_{B/k}.
\end{equation}
One can also obtain this directly from the universal properties:
the displayed inverse Chow operations give inverse homomorphisms
of the three universal regular representatives.

Let $b:B_K\hookrightarrow Y$. Under
\eqref{eq:J6P-blowup}, $\delta$ is the sum of $j^*$ and the
homomorphism induced by $b_*$ on algebraically trivial divisors.
There is no excess term here: $E$ and the smooth generic fibre
meet in the Cartier divisor $B_K$, so refined base change gives
$\delta e_*q^*=b_*$.

Choose a smooth member $Y_0/k$ of the pencil; one exists since
the smooth locus of the proper pencil contains its generic point.
Applying weak Lefschetz to $X\supset Y_0\supset B$ shows that
restriction
\[
 a:\operatorname{Pic}^{\rm var}_{X/k}
       \longrightarrow\operatorname{Pic}^{\rm var}_{B/k}
\]
is an isogeny. Indeed, $n\geq4$ makes both restrictions
isomorphisms on $H^1(-,\mathbb Q_\ell(1))$ and hence on the
rational Tate modules of these Picard varieties. Write $H$ for
the chosen hyperplane class. The projection formula gives,
for every algebraically trivial line bundle $L$ on $X$,
\[
 b_*(L|_{B_K})=H|_Y\cdot L|_Y=j^*(H\cdot L).
\]
Thus, as homomorphisms of abelian varieties,
$b_*a_K=j^*\beta_K$, where
$\beta:\operatorname{Pic}^{\rm var}_{X/k}\to
\operatorname{Ab}^2_{X/k}$ is induced by $L\mapsto H\cdot L$.
Since $a$ is surjective, the exceptional summand has image in
$\operatorname{im}(j^*)$. Therefore
$\operatorname{im}(\delta)=\operatorname{im}(j^*)$.

Equation \eqref{eq:J6P-spread} and the surjectivity of $[r]_T$
now imply $\operatorname{im}(\tau)\subseteq\operatorname{im}(j^*)$.
Conversely, the map $j^*$ has constant source; the universal
property of $T$ factors it through $\tau$. This proves
\eqref{eq:J6P-trace}.

\emph{Step 3: why the trace fills the stable-range target.}
Assume $n\geq5$. Choose a prime $\ell$ different from the
characteristic and coprime to the integer of a miniversal cycle
for $A$. By \cite[Proposition 8.2, p.~31]{DGJ_ACMVFunctor}, there
is a Galois-equivariant injection
\[
 T_\ell A\hookrightarrow H^3(Y_{K^{\rm sep}},\mathbb Z_\ell(2))_{
 \rm torsion\,free}.
\]
Weak Lefschetz identifies its rational target with
$H^3(X_{K^{\rm sep}},\mathbb Q_\ell(2))$. This Galois
representation is constant because $X$ is defined over the
algebraically closed field $k$. Thus $T_\ell A$ has trivial
Galois action. All $\ell$-power torsion points of $A$ are
$K$-rational, and their union is a divisible group.

By Lang--N\'eron,
$A(K)/\tau(T(k))$ is finitely generated
\cite[Theorem 7.1, p.~23]{DGJ_Conrad2006}. A divisible group
has zero image in a finitely generated abelian group, so
$A(K)[\ell^\infty]\subseteq\tau(T(k))$. Prime-to-characteristic
power torsion is Zariski dense in an abelian variety. For
completeness, if a proper abelian subvariety contained all this
torsion, its positive-dimensional abelian quotient would have
no $\ell$-power torsion; this is impossible since torsion lifts
through a surjective homomorphism of abelian varieties over an
algebraic closure. Hence $\operatorname{im}(\tau)=A$, as claimed.

\emph{The finite-kernel boundary.}
For $n\geq4$, the same $j^*$ has finite kernel if the universal
map for $Y$ is injective on $\ell$-primary cycle torsion for a
prime $\ell$ avoiding the characteristic and the integer of a
chosen miniversal cycle for $X$. Here is the full implication.
If the reduced identity component $C$ of its kernel were
positive-dimensional, restrict a miniversal cycle $\Gamma_X$
to $C$. The map
\[
 T_\ell C\longrightarrow T_\ell A^2(X_{K^{\rm sep}})
 \longrightarrow H^3(X_{K^{\rm sep}},\mathbb Z_\ell(2))_{
 \rm torsion\,free}
\]
is injective by Proposition 8.2. Its pullback cycle family on
$Y$ is killed by $\Phi_Y$ since the representative map is
$[r]j^*|_C=0$. The stated torsion injectivity then makes its
Tate module of cycles zero. Functoriality of the Bloch map
\cite[Section 8.1, pp.~30--31]{DGJ_ACMVFunctor} makes the
displayed cohomology classes restrict to zero on $Y$. But weak
Lefschetz is injective on $H^3(X)\to H^3(Y)$ for every $n\geq4$,
a contradiction after tensoring with $\mathbb Q_\ell$.

This torsion hypothesis is known in characteristic zero by
Murre's theorem, as recorded in
\cite[Remark 8.6, p.~33]{DGJ_ACMVFunctor}. It therefore gives,
in characteristic zero and with the smooth-axis hypothesis,
an isogeny in the stable range and an isogeny onto the trace
image when $n=4$. Proposition 8.2 alone does not remove the
hypothesis: its injection is constructed from a miniversal
cycle and does not assert that $\Phi_Y$ is injective on torsion.
The same remark explicitly leaves that injectivity unsettled
in general positive characteristic.

\paragraph{J6Z (The highest nontrivial target codimension).}
Let $X/k$ be smooth projective of dimension $n\geq3$ and let
$j:Y\hookrightarrow X$ be any smooth hyperplane section. Then
\[
 j^*:P^{n-1}_{\rm inc}(X)\longrightarrow
       P^{n-1}_{\rm inc}(Y)=\operatorname{Alb}(Y)
\]
is surjective in every characteristic. The analogous assertion
holds for universal regular representatives whenever the source
exists; in particular it holds for
$\operatorname{Ab}^2(X)\to\operatorname{Alb}(Y)$ on a threefold.

To prove it, restriction
$\operatorname{Pic}^{\rm var}(X)\to\operatorname{Pic}^{\rm var}(Y)$
is an isogeny by weak Lefschetz in degree one. Intersection by
$H_X^{n-2}$ defines a homomorphism from the former Picard
variety to $P^{n-1}_{\rm inc}(X)$, and its restriction to $Y$
is intersection by $H_Y^{n-2}$ followed by the Albanese map.
The induced map
\[
 \operatorname{Pic}^{\rm var}(Y)\longrightarrow\operatorname{Alb}(Y)
\]
is an isogeny: on rational Tate modules, the codimension-one
and zero-cycle Bloch isomorphisms identify it with
\[
 L_Y^{n-2}:H^1(Y,\mathbb Q_\ell(1))
       \xrightarrow{\ \sim\ }H^{2n-3}(Y,\mathbb Q_\ell(n-1)),
\]
which is hard Lefschetz. See
\cite[Section 8.1, pp.~30--31]{DGJ_ACMVFunctor} for those
Bloch isomorphisms and
\cite[Theorem 4.1.1 and (4.1.6), pp.~217--219]{DGJ_Deligne1980}
for hard and weak Lefschetz. A surjective homomorphism thus
factors through $j^*$, proving surjectivity. In this top target
codimension, the common incidence quotient is the Albanese:
the Poincar\'e family of divisors on $Y$ detects exactly its
kernel. This argument uses neither an inverse Lefschetz
correspondence nor the torsion standard assumption.

\emph{Exact scope.}
The fixed-part result proves the modern codimension-two
interpretation of the middle clause for pencils with smooth
axis. It does not identify that interpretation with the
manuscript's undefined equivalence relation, cover singular
pencil axes, prove positive-characteristic finite kernel, or
settle the higher-codimension stable and middle assertions.

% END J6P

% BEGIN J6C
\subsection{A cubic pencil refuting the literal middle fixed-part clause}
\paragraph{J6C (Refuted by an explicit cubic pencil).}
At 44:4, Theorem~6 says, under the conditions of Theorem~5,
``si $2i=n$ le facteur fixe de $P^i(Y)$ est identique à
l'image de $P^i(X)$''. The inserted condition on $Y$ in
Theorem~5 is ``section hyperplane générique d'un pinceau''
whose fibres are geometrically integral. No smoothness of the
pencil axis is imposed in either supplied transcription.
The word ``générique'' modifies the hyperplane section of the
pencil; the pencil itself is not required to be general.
The following example meets those literal hypotheses and has
zero source but a positive-dimensional fixed target factor.
The cycle hypotheses in Theorem~5 can, if needed, be met by
taking $Z=Z'=0$ with $i=j=2$.

\emph{Existing-source check.}
The underlying elliptic-cone geometry is classical:
\cite[Lemma 8.1, pp.~314--315]{DGJ_ClemensGriffiths1972}
describes the cone of lines through an Eckardt point, and
\cite[Proposition 10, p.~384; Lemma 29, p.~398]{DGJ_Roulleau2009}
studies the associated elliptic curve and homomorphisms on the
complex Fano surface. The explicit pencil and positive-characteristic
calculation below are supplied with a complete proof; no priority
claim is made for this construction.

Let $k=\overline{\mathbb F}_5$, and let $X\subset\mathbb P^5_k$
be the cubic fourfold
\begin{equation}\label{eq:J6C-fourfold}
 x_0^3+x_1^3+x_2^3+x_3^2x_4+x_4^3+x_5^3=0.
\end{equation}
Consider the hyperplane pencil $x_5=t x_4$, with parameter
$t\in\mathbb A^1_k$, and the member $x_4=0$ at infinity.
Its axis is the integral surface
\[
 B:\quad x_4=x_5=0,\qquad x_0^3+x_1^3+x_2^3=0.
\]
It is the cone over the smooth plane cubic
\[
 E:\quad u_0^3+u_1^3+u_2^3=0\quad\text{in }\mathbb P^2_k,
\]
with vertex $p=[0:0:0:1:0:0]$. Choose
$e_0=[1:-1:0]$ on $E$ as origin. Let $K=k(t)$, and let
$Y/K$ be the generic pencil member. The bundle
$\mathcal O_X(1)$ is the very ample bundle of this displayed
embedding. The morphisms below are all defined over $K$.

\emph{Smooth total variety and integral pencil members.}
The gradient of \eqref{eq:J6C-fourfold} is
\[
 (3x_0^2,3x_1^2,3x_2^2,2x_3x_4,
   x_3^2+3x_4^2,3x_5^2).
\]
A common zero forces $x_0=x_1=x_2=x_5=0$, then
$x_3x_4=0$ and $x_3^2+3x_4^2=0$, and finally
$x_3=x_4=0$. Thus $X$ is smooth. For a finite parameter,
the member is the cubic in $\mathbb P^4$ with equation
\begin{equation}\label{eq:J6C-member}
 x_0^3+x_1^3+x_2^3+x_3^2x_4+a x_4^3=0,
 \qquad a=1+t^3.
\end{equation}
If $a\ne0$, the same derivative calculation shows that it
is smooth. If $a=0$, its singular locus has the single
geometric point $[0:0:0:0:1]$. The member at infinity is
the cone, with vertex $p$, over the smooth Fermat cubic
surface $x_0^3+x_1^3+x_2^3+x_5^3=0$; its singular locus
also has just one point.

Every member is geometrically integral. Indeed, over an
algebraically closed field a reducible hypersurface in
$\mathbb P^4$ has two components meeting in dimension at
least two, and their intersection is singular for the
product equation. A repeated component is singular along
a three-dimensional subset. Both possibilities contradict
the computed empty or zero-dimensional singular locus.
This also proves geometric integrality of the generic
member after extending $K$ to an algebraic closure. The
three finite singular parameters are distinct, since
$1+t^3$ is relatively prime to $3t^2$ in characteristic five.

The pencil is the incidence family in $X\times\mathbb P^1$.
Because $x_4,x_5$ form a regular sequence along the
codimension-two axis, the two-generated ideal is of linear
type \cite[Proposition IV-25, pp.~172--173]{DGJ_EisenbudHarris2000};
its blowup is that incidence family. Thus its scheme
fibres are exactly the displayed hyperplane sections,
with no additional fibre components. Equivalently, each
local incidence equation restricts to a nonzerodivisor on
each fibre of $X\times\mathbb P^1$, so this is a flat
hyperplane family. The singular axis is the sole missing
hypothesis relative to the positive theorem J6P.

\emph{The source representative vanishes.}
Weak Lefschetz for the smooth ample hypersurface
$X\subset\mathbb P^5$ gives
\[
 H^3(X,\mathbb Q_\ell)=H^3(\mathbb P^5,\mathbb Q_\ell)=0
 \qquad(\ell\ne5).
\]
Consequently
\begin{equation}\label{eq:J6C-zero-source}
 \operatorname{Ab}^2(X)=0,\qquad P^2_{\rm inc}(X)=0.
\end{equation}
For the first conclusion use the miniversal-cycle injection
\cite[Proposition 8.2, p.~31]{DGJ_ACMVFunctor}; for the
second use Saito's incidence bound
\cite[Proposition 3.4, pp.~109--110]{DGJ_Saito1979}.
Weak Lefschetz is cited in
\cite[(4.1.6), pp.~218--219]{DGJ_Deligne1980}.

\emph{A constant elliptic factor in the target.}
Blow up the $K$-point $p$ of the smooth threefold $Y$:
\[
 \pi:\widehat Y=\operatorname{Bl}_pY\longrightarrow Y,
 \qquad D\simeq\mathbb P^2_K.
\]
Let $i:S\hookrightarrow\widehat Y$ be the strict transform
of $B_K$. This is the smooth ruled resolution of the
elliptic cone, with projection $q:S\to E_K$. Concretely,
its fibre over $e=[u_0:u_1:u_2]$ is the strict transform
of the ruling line
\[
 [\lambda u_0:\lambda u_1:\lambda u_2:\mu:0:0]
 \quad\text{in }B_K.
\]
Thus $S$ is a $\mathbb P^1$-bundle over $E_K$, and every
ruling fibre $\ell$ satisfies
$\pi^*H\cdot\ell=1$ and $D\cdot\ell=1$.

The multiplicity of $B_K$ as a divisor of $Y$ at $p$ is
three. To check it, use $x_3=1$ and put $v=x_4$; equation
\eqref{eq:J6C-member} becomes
\[
 v+a v^3+u_0^3+u_1^3+u_2^3=0.
\]
The three $u_j$ are regular parameters at $p$, and
$v(1+a v^2)=-(u_0^3+u_1^3+u_2^3)$, with the first factor
in parentheses a unit. Hence $v$ has order three. As a
Cartier divisor on $\widehat Y$, therefore,
\[
 S=\pi^*H-3D,
 \qquad
 \deg(N_{S/\widehat Y}|_\ell)=1-3=-2.
\]

Put $J=\operatorname{Pic}^0(E)$, a constant elliptic
curve over $k$. The algebraic-correspondence operations
\[
 i_*q^*:A^1(E_K)\longrightarrow A^2(\widehat Y),\qquad
 q_*i^*:A^2(\widehat Y)\longrightarrow A^1(E_K)
\]
preserve algebraic equivalence and regularity. Here $i^*$
is the refined pullback for the Cartier embedding $i$;
$q_*$ sends zero-cycles on the surface $S$ to divisor
classes on $E_K$. They induce homomorphisms of abelian
varieties
\[
 \widehat u:J_K\longrightarrow\operatorname{Ab}^2(\widehat Y),
 \qquad
 \widehat v:\operatorname{Ab}^2(\widehat Y)\longrightarrow J_K.
\]
The Cartier self-intersection formula
\cite[Definition 42.29.1(1)(b)]{DGJ_StacksGysin}
and the projection formula
\cite[Lemma 42.26.4]{DGJ_StacksProjection} give,
for every $L\in\operatorname{Pic}^0(E_K)$,
\begin{align*}
 q_*i^*i_*q^*L
 &=q_*\bigl(c_1(N_{S/\widehat Y})\cdot q^*L\bigr)\\
 &=\bigl(q_*c_1(N_{S/\widehat Y})\bigr)\cdot L=-2L.
\end{align*}
Thus
\begin{equation}\label{eq:J6C-minus-two}
 \widehat v\widehat u=[-2]_{J_K}.
\end{equation}
This sign is an equality of divisor-class homomorphisms
on $J$, not an unstated choice of sign for an elliptic
principal polarization.

The point-blowup formula gives
\[
 \operatorname{CH}^2(\widehat Y)
    =\pi^*\operatorname{CH}^2(Y)\oplus\mathbb Z[\text{line in }D].
\]
The exceptional integer is constant under algebraic
equivalence, so $\pi^*$ and $\pi_*$ are inverse on $A^2$.
They also preserve regular families and therefore induce
inverse isomorphisms of the universal regular
representatives. This is the $r=3$ case of the blowup
formula \cite[Theorem 2.8(3)]{DGJ_Jiang2023}.
Define
\[
 u=\pi_*\widehat u:J_K\longrightarrow\operatorname{Ab}^2(Y),
 \qquad
 v=\widehat v\pi^*:\operatorname{Ab}^2(Y)\longrightarrow J_K.
\]
Equations above give $vu=[-2]_{J_K}$. In characteristic
five this map is an isogeny of degree four. Hence $u$ has
finite kernel and a one-dimensional abelian image. Since
its source is constant and $u$ is defined over $K$, the
universal property of the $K/k$-trace puts this image in
the fixed part of $\operatorname{Ab}^2(Y)$. That fixed
part is therefore positive-dimensional, contradicting
\eqref{eq:J6C-zero-source}.

The factor is detected by incidence equivalence as well:
$q_*i^*\pi^*$ is an actual correspondence from $Y$ to the
curve $E_K$, taking codimension-two classes to divisor
classes. It is one of the Picard tests defining the
incidence quotient. Thus the constant family $i_*q^*$
cannot vanish in that quotient, since composing with
this test is $[-2]$. In particular the counterexample
does not depend on the disputed integral autoduality
of a universal representative.

\emph{Conclusion and remaining scope.}
For $n=4,i=2$, a generic hyperplane section of this
integral-fibre pencil has a nonzero fixed elliptic
factor, although the ambient fourfold has zero
codimension-two representative. The literal middle
clause of Theorem~6 is therefore false. This does not
refute J6P, whose smooth-axis hypothesis is essential,
or settle the separate stable- and high-range assertions
in all codimensions. The exact gradient ideals,
singular supports, origin, parameter roots, and normal
degree are checked by \texttt{check\_cubic\_trace.py};
the constant-factor conclusion is the correspondence
proof above, not an output of the finite computation.

% END J6C

% BEGIN JWalker
\subsection{The remaining regular-representability gate}
The August 5, 2026 revision of
\cite[Introduction, p.~1]{DGJ_ACMVTwist2026} still distinguishes
universal algebraic representatives in codimensions $1,2,\dim X$
from merely regular algebraic intermediate-Jacobian maps in
the other codimensions. This is a precise current source check,
not an inference of open status from an unsuccessful search.
The final published version of Kahn's concrete fourfold
expectation is now also opened: Conjecture~3.1, p.~847
\cite{DGJ_Kahn2021}. It corresponds to Conjecture~1, p.~5 in
the six-page author version cited above.

For complex varieties one should additionally distinguish
universality among all regular homomorphisms from universality
among isogeny lifts. Write
\[
 \psi:A^i(X)\longrightarrow J_a^{2i-1}(X)(\mathbb C),\qquad
 \omega:A^i(X)\longrightarrow J_W^{2i-1}(X)(\mathbb C)
\]
for the usual and Walker Abel--Jacobi maps, and
$\alpha:J_W^{2i-1}(X)\to J_a^{2i-1}(X)$ for their isogeny.
By \cite[Theorem 2.3, p.~12]{DGJ_ACMVWalker}, $\omega$ is
initial among regular lifts of $\psi$ along isogenies. This
theorem does not say that $\omega$ is universal among all
regular homomorphisms.

Here is an exact elementary test for that stronger assertion.
For every regular homomorphism
$\phi:A^i(X)\to B(\mathbb C)$, let
\[
 C_\phi=\operatorname{im}(\psi,\phi)
       \subset J_a^{2i-1}(X)\times B.
\]
The image is the group of complex points of an abelian
subvariety, by the image theorem for regular homomorphisms
\cite[Lemma 4.7, pp.~18--19]{DGJ_ACMVFunctor}.
Then $\omega$ is universal among all regular homomorphisms
if and only if
\begin{equation}\label{eq:JWalker-gate}
 \dim C_\phi=\dim J_a^{2i-1}(X)
 \quad\hbox{for every regular }\phi.
\end{equation}
Indeed, if \eqref{eq:JWalker-gate} holds, the surjective
projection $C_\phi\to J_a^{2i-1}(X)$ is an isogeny. Initiality
among isogeny lifts gives a homomorphism
$J_W^{2i-1}(X)\to C_\phi$ inducing $(\psi,\phi)$; its second
projection factors $\phi$ through $\omega$. This factorization
is unique because $\omega$ is surjective. Conversely, if
$\omega$ is universal, $(\psi,\phi)$ factors through $J_W$,
so its image has dimension at most $\dim J_W=\dim J_a$;
the surjective projection to $J_a$ gives the reverse inequality.

The unresolved issue in this route is therefore a possible
positive-dimensional kernel of $C_\phi\to J_a$ for a regular
test $\phi$. The divisible-kernel statement in Theorem~2.3
does not rule it out: divisible groups can have nonzero maps
to groups of points of abelian varieties. Likewise, infinite
torsion in $A^3$ does not by itself construct regular targets
of unbounded dimension. The test above is not asserted for
Kahn's fourfold, and neither its proof nor its failure has
been obtained here. This leaves the general statement in
Theorems~1--3 explicitly unresolved.

% END JWalker
